\section{Controllers}

\slidebackground{backgrounds/wallpaper.jpg}

\begin{frame}[fragile]{What Is a Controller?}
	\begin{itemize}
		\item A \textbf{controller} is a Python class that handles HTTP requests.
		\item In Odoo it inherits from \texttt{http.Controller}.
		\item Controllers define routes that return HTML, JSON, files, or redirects.
\begin{block}{Method Syntax}
\begin{lstlisting}
from odoo import http
from odoo.http import request

class StudentController(http.Controller):
	pass
\end{lstlisting}
\end{block}
	\end{itemize}
\end{frame}

\begin{frame}[fragile]{Controller File Placement}
\begin{lstlisting}
school_manager/
  controllers/
	__init__.py
	student_portal.py
  __init__.py
\end{lstlisting}
\begin{lstlisting}
# school_manager/__init__.py
from . import controllers

# school_manager/controllers/__init__.py
from . import student_portal
\end{lstlisting}
	\begin{itemize}
		\item Controllers must be imported or Odoo will not register the routes.
	\end{itemize}
\end{frame}

\begin{frame}[fragile]{request Object}
	\begin{itemize}
		\item Inside controllers, \texttt{request} is the current HTTP request.
\begin{lstlisting}
request.env          # environment for current user/db
request.httprequest  # underlying Werkzeug request
request.session      # session data
request.params       # combined request parameters
\end{lstlisting}
		\item Most business work uses \texttt{request.env['model.name']}.
	\end{itemize}
\end{frame}

\begin{frame}[fragile]{Simple HTML Controller}
\begin{lstlisting}
class StudentController(http.Controller):

	@http.route('/school/students', type='http', auth='public', website=True)
	def student_list(self, **kwargs):
		students = request.env['student.student'].sudo().search([])
		return request.render('school_manager.student_list_page', {
			'students': students,
		})
\end{lstlisting}
	\begin{itemize}
		\item \texttt{request.render()} renders a QWeb website template.
	\end{itemize}
\end{frame}

\begin{frame}{Controller Best Practices}
	\begin{itemize}
		\item Keep controllers thin; put business logic in models.
		\item Choose \texttt{auth} carefully (\texttt{public}, \texttt{user}, \texttt{none}).
		\item Validate input before writing to the database.
		\item Use \texttt{sudo()} only when the route intentionally needs elevated access.
	\end{itemize}
\end{frame}
